%==============================================================================
\section{Frontend: Live Dashboard}
%==============================================================================

The dashboard is a single self-contained page,
\texttt{ecg/templates/ecg.html}, served at the site root. It pulls Chart.js
4.4.1 from a CDN and draws the heart waveform on a dark background, refreshed in
real time as cleaned batches arrive over the WebSocket.

\subsection{Rolling buffer and chart}
The page keeps a fixed-length \emph{rolling buffer} of 1000 samples (about four
seconds of signal at 250\,Hz). Each incoming batch is appended to the right while
the oldest samples drop off the left, producing the familiar scrolling-trace
effect of a bedside monitor. The chart's $y$-axis is pinned to the ADC range
$0$--$4095$ so the waveform's scale is stable, and animation is disabled so
updates are instantaneous.

\begin{table}[h]
  \centering\small
  \begin{tabular}{@{}ll@{}}
    \toprule
    \textbf{Setting} & \textbf{Value} \\
    \midrule
    Charting library   & Chart.js 4.4.1 (CDN) \\
    Chart type         & line, no point markers, slight tension (0.2) \\
    Buffer size        & 1000 samples ($\approx$4\,s window) \\
    $y$-axis range     & 0--4095 (fixed, matches ADC) \\
    $x$-axis           & linear sample index, hidden \\
    Animation          & disabled (real-time redraw) \\
    Trace colour       & lime on near-black background \\
    \bottomrule
  \end{tabular}
  \caption{Dashboard chart configuration.}
  \label{tab:chart}
\end{table}

\subsection{Receiving data}
The page opens a WebSocket back to the same host it was served from, so no
hard-coded server address is needed in the browser. Each message is a
\texttt{\{"filtered":[...]\}} array; the handler shifts the buffer and redraws.

\begin{lstlisting}[style=js,caption={Dashboard WebSocket handler and buffer update (abridged).}]
const BUFFER_SIZE = 1000;
let ecgBuffer = Array(BUFFER_SIZE).fill(0);

const ws = new WebSocket("ws://" + window.location.host + "/ws/ecg/");

ws.onmessage = (event) => {
  const msg = JSON.parse(event.data);
  const incoming = msg.filtered || msg.ecg;          // cleaned samples
  if (Array.isArray(incoming)) {
    // drop oldest, append newest -> scrolling trace
    ecgBuffer = ecgBuffer.slice(incoming.length).concat(incoming);
    ecgChart.data.datasets[0].data = ecgBuffer;
    ecgChart.update();
  }
};
\end{lstlisting}

Because the dashboard subscribes to the very same \texttt{ecg\_group} the server
broadcasts to, opening the page in several browsers shows the same live trace in
all of them simultaneously -- this is what makes the monitor ``remote'': any
machine that can reach the server can watch the patient's ECG.
